\section*{2 Preliminaries}
We introduce two space-bounded models of computation. Our first model is a variation of branching programs, which are the standard syntactic (i.e. non-uniform) notion of space-bounded computation (see [7]).
\paragraph{Definition5 ($m$-catalytic branching program [8]).} Let $n\in\mathbb N$ and let $f:\{0,1\}^n\to\{0,1\}$ be an arbitrary function. An $m$-catalytic branching program is a directed acyclic graph $G$ with the following properties:
\begin{itemize}
\item There are $m$ source nodes and $2m$ sink nodes.
\item Every non-sink node is labeled with an input variable $x_i$ for $i\in[n]$, and has two outgoing edges labeled0 and1.
\item For every source node $v$ there is one sink node labeled with $(v,0)$ and one with $(v,1)$.
\end{itemize}
We say that $G$ computes $f$ if for every $x\in\{0,1\}^n$ and source node $v$, the path defined by starting at $v$ and following the edges labeled by the value of the $x_i$ labeling each node ends at the sink labeled by $(v,f(x))$. The size of $G$ is the number of non-sink nodes\footnote{This is somewhat non-standard, but when talking about layered branching programs this simplifies things by defining the length as the number of times we read variables, which will in turn be connected to the number of instructions in our register program model. This choice does not affect the asymptotics of any results.} in $G$.
We also consider $m$-catalytic branching programs with standard restrictions:
\begin{itemize}
\item layered branching programs: for some $\ell\in\mathbb N$, there exists a function $\sigma:G\to[\ell+1]$ such that for all $u\in G$, the outgoing edges of $u$ go to nodes $v,w$ such that $\sigma(v)=\sigma(w)=\sigma(u)+1$; we call the set of nodes $\{v\in\sigma^{-1}(j)\}$ the $j$th layer. Furthermore for each $j\in[\ell]$ there exists an input variable $x_{j_i}$ which is the variable labeling every $v\in\sigma^{-1}(j)$.\footnote{By construction layer $\ell+1$ will contain exactly the sink nodes of $G$. See the previous footnote for an explanation of this somewhat non-standard convention.} The width of $G$ is $\max_{j\in[\ell]}|\sigma^{-1}(j)|$ and the length of $G$ is $\ell$. Note that the size of $G$ is at most the product of the length and width of $G$.
\item read-$k$ branching programs: for any start node $v$ and any input $x$, each variable $x_i$ is read at most $k$ times during the computation starting at $v$. We note that for layered branching programs, this corresponds to each variable $x_j$ labeling at most $k$ layers.
\item permutation branching programs: consider a layered branching program of width $2m$ with $2m$ source nodes $S=\{s_1\ldots s_{2m}\}$ and $2m$ sink nodes $T=\{t_1\ldots t_{2m}\}$, and let $P(x):[2m]\to[2m]$ be the function such that $P(x)(i)=i'$ where the computation path starting from $s_i$ on input $x$ goes to $t_{i'}$. Then there exists a permutation $\sigma^*:[2m]\to[2m]$ such that $P(x)$ is the identity function when $f(x)=0$ and $P(x)=\sigma^*$ when $f(x)=1$. This is not an $m$-catalytic branching program in the strict sense, and we will address this model more precisely in Section4.
\end{itemize}
Our second model comes from a line of work starting with [3], more recently fleshed out in [4] and used in many follow-up works on catalytic computation [5,6].
\paragraph{Definition6 (Transparent register program).} Let $R$ be a ring and $f\in\{0,1\}^n\to\{0,1\}$ a function. A register program $P$ is defined by a set of registers $\mathcal S=\{R_1\ldots R_s,R_{\mathrm{out}}\}$, each storing a value in $R$, plus an ordered list of $t$ instructions where for every $j\in[t]$ the $j$th instruction is $R\leftarrow R+p_j(x_i,\mathcal S\setminus\{R\})$ for some $i\in[n]$, $R\in\mathcal S$, and polynomial $p_j\in R[x,y_0\ldots y_s]$.
We say that $P$ computes $f$ if for every $x\in\{0,1\}^n$, after initializing $R=0$ for all registers $R$ and then executing all instructions in order, the value stored in $R_{\mathrm{out}}$ is $f(x)$. We say that $P$ transparently computes $f$ if instead of initializing all $R$ to0, each $R$ begins in an arbitrary initial state $\tau_i$, and at the end of the program we have $R_{\mathrm{out}}=\tau^{\mathrm{out}}+f(x)$ and $R_i=\tau_i$ for all $i$. The size of $P$ is the number of registers $s+1$ and the time of $P$ is the number of instructions $t$.
We use register programs as an abstraction for understanding $m$-catalytic branching programs, which will be the model our main results will refer to. The two models are connected through the following observation (see Figure2 of the original article for an example of our construction).
\paragraph{Observation7.} Let $f_n:\{0,1\}^n\to\{0,1\}$ be a family of functions and let $P_n$ be a family of register programs over a family of rings $R_n$ with size $s(n)+1$ and time $t(n)$ each transparently computing $f_n$. Then $f_n$ can be computed by a family of $m$-catalytic branching programs of width $|R_n|m$ and length $t(n)$, where $m=|R_n|^{s(n)}$.
\begin{proof}
We will define a branching program with width $|R_n|^{s(n)+1}$ and length $t(n)$. Each layer will represent a stage in the register program and each node in a layer will represent a setting to the registers at that time. Since each register program step only requires us to read one input variable, at layer $k$ we query the variable associated with step $k$ in the register program and create edges from each node in layer $k$ to the nodes at layer $k+1$ representing the state of our memory after the step has completed. We label each input node as $\tau$ for some distinct initial configuration $\tau=(\tau_1\ldots\tau_{s(n)})$ to all registers except $R_{\mathrm{out}}$, and we treat $R_{\mathrm{out}}$ as being initialized to0. Then by the fact that $P$ transparently computes $f$, starting at node $\tau$ we are guaranteed to reach a node $(\tau,f(x))$. Since in each layer we have a node for every setting of the $s(n)+1$ registers, the width of our branching program is $|R_n|^{s(n)+1}$, and since each non-output layer corresponds to a unique instruction from our register program, the length of our branching program is $t(n)$.
\end{proof}
\paragraph{Note8.} In our computation we often include instructions of the form $R\leftarrow R+p_j(R/\{R\})$, i.e. instructions that do not require reading a variable $x_i$. By observing the above proof, it is clear that such instructions do not add any length to the branching program, as they can be computed at the same time as an adjacent instruction.
